\section{Common-information stability of a positive Gaussian correction}
\label{sec:common-information-stability}

I separate the mathematics of a common shift from the finite arithmetic
which implements it. A common shift is a perturbation of both endpoint
arguments by the same amount; it need not be probabilistically independent
of the outer financial coordinate or of the gap. The essential condition
for Gaussian smoothing is that, after conditioning on the outer labels
and any fixed strike label, the shift and gap are independent of the
inner normal label. A label-dependent scaling remainder does not satisfy
this condition and receives a separate allowance.

Let $\mu$ be a positive finite Borel measure of mass one, let $\lambda>0$,
and define
\begin{equation}
 P_{\lambda,\mu}(m)=\int_{\mathbb R}
       \Phi((m-a)/\lambda)\,\mu(da),\qquad
 R_{\lambda,\mu}(m,\Delta)=P_{\lambda,\mu}(m+\Delta)
                              -P_{\lambda,\mu}(m).
 \label{eq:common-information-definition}
\end{equation}
This includes a fixed strike, a finite positive strike mixture and a
bounded increasing payoff normalised to total variation one. In the last
case $\mu=df$, where $f(-\infty)=0$ and $f(+\infty)=1$; Tonelli's theorem
identifies $P_{\lambda,\mu}(m)$ with $\mathbb E f(m+\lambda N)$.
These identities and Gaussian convolution are established mechanisms.
The calculation below concerns their shared perturbations and their
specific guarded pricing implementations; it is not a claim of priority
for convolution or Lipschitz estimates.

Choose $\nu_0\geq\sup_x\phi(x)$ and
$\nu_1\geq\sup_x|\phi'(x)|$, and put
\begin{equation}
 c_0=\frac{\nu_0}{\lambda},\qquad
 c_1=\frac{\nu_1}{\lambda^2},\qquad
 C(d)=\min\{c_0,c_1d\}\quad(d\geq0).
 \label{eq:common-information-constants}
\end{equation}

\begin{theorem}[Common-shift and gap stability]
\label{thm:common-information-stability}
For finite $m,\widehat m,\Delta,\widehat\Delta$ with $\Delta\geq0$, let
$|\widehat m-m|\leq E_m$ and
$|\widehat\Delta-\Delta|\leq e_\Delta$. The represented gap may be
negative. Then
\begin{equation}
 |R_{\lambda,\mu}(\widehat m,\widehat\Delta)
                       -R_{\lambda,\mu}(m,\Delta)|
 \leq C(\Delta)E_m+c_0e_\Delta.
 \label{eq:common-information-stability}
\end{equation}
Consequently a uniform true-gap bound $0\leq\Delta\leq\Delta_*$ permits
the coefficient $C(\Delta_*)$. If $\widehat\Delta\geq0$ as well, the
coefficient in \eqref{eq:common-information-stability} can be replaced
by $C(\min\{\Delta,\widehat\Delta\})$.
Let $\Pi_{[0,B]}$ be metric projection onto $[0,B]$, where $B>0$ and
$R_{\lambda,\mu}(m,\Delta)\in[0,B]$. The same upper bound applies after
projecting the represented difference onto $[0,B]$. Projection need not
be inserted into the implemented gap register.
\end{theorem}
\begin{proof}
The first and second derivatives of the Gaussian kernel are bounded
uniformly in $a$ by $c_0$ and $c_1$. Difference quotients of the kernel
are dominated by these constants by the mean value theorem. Since
$\mu$ is finite, dominated convergence justifies differentiation under
its integral twice, giving
\begin{align}
 P_{\lambda,\mu}'(m)&=\int\lambda^{-1}
                 \phi((m-a)/\lambda)\,\mu(da),\nonumber\\
 P_{\lambda,\mu}''(m)&=\int\lambda^{-2}
                 \phi'((m-a)/\lambda)\,\mu(da),\qquad
 0\leq P_{\lambda,\mu}'\leq c_0,\quad
 |P_{\lambda,\mu}''|\leq c_1.
 \label{eq:common-information-derivatives}
\end{align}
The difference of two numbers in $[0,c_0]$ has absolute value at most
$c_0$. The fundamental theorem of calculus also bounds the difference
of the derivatives at arguments separated by a nonnegative distance $d$
by $c_1d$. Thus
\begin{equation}
 |\partial_m R_{\lambda,\mu}(m,d)|
   =|P_{\lambda,\mu}'(m+d)-P_{\lambda,\mu}'(m)|\leq C(d),\qquad
 |\partial_dR_{\lambda,\mu}(m,d)|\leq c_0.
 \label{eq:common-information-partial-bounds}
\end{equation}
For the gap partial derivative the bound holds for all real $d$.
Split the difference exactly as
\begin{align}
 R(\widehat m,\widehat\Delta)-R(m,\Delta)
 &=R(\widehat m,\widehat\Delta)-R(\widehat m,\Delta)
       +R(\widehat m,\Delta)-R(m,\Delta).
 \label{eq:common-information-path}
\end{align}
Integrating the gap derivative along the first segment gives
$c_0|\widehat\Delta-\Delta|$, including a segment crossing zero.
Integrating the common-shift derivative along the second segment, whose
gap is the true nonnegative $\Delta$, gives
$C(\Delta)|\widehat m-m|$. This proves the stated inequality. If the
represented gap is nonnegative, reversing the order of the two segments
gives $C(\widehat\Delta)|\widehat m-m|+c_0|\widehat\Delta-\Delta|$;
the smaller valid bound gives the refinement. Since $C$ is nondecreasing,
a valid upper gap bound may replace the gap. Finally, metric projection
is $1$-Lipschitz and fixes the true correction. Applying it to the
represented difference cannot increase its distance from that correction.
No cancellation of two distinct scalar-table approximation errors is
asserted.
\end{proof}

\begin{lemma}[Common inner perturbations of capped payoffs]
\label{lem:common-information-inner-payoff}
For the normalised cap
$f(x)=H^{-1}\min\{K(e^x-1)_+,H\}$, put
$\ell=\log(1+H/K)$ and $L_f=(K+H)/H$. For every $d\geq0$, the map
$g_d(x)=f(x+d)-f(x)$ is globally $L_f$-Lipschitz. For any real
perturbation $\zeta$, which may depend on the inner label,
\begin{equation}
 |g_d(x+\zeta)-g_d(x)|\leq L_f|\zeta|.
 \label{eq:common-information-inner-payoff}
\end{equation}
\end{lemma}
\begin{proof}
The payoff is absolutely continuous and has almost-everywhere derivative
$f'(x)=(K/H)e^x\mathbf1_{(0,\ell)}(x)$, which lies in $[0,L_f]$.
The difference $g_d$ is absolutely continuous on every finite interval
and its almost-everywhere derivative is $f'(x+d)-f'(x)$.
Its absolute value is at most $L_f$, since both terms belong to the same
interval $[0,L_f]$. Integrating that derivative between $x$ and $x+\zeta$
proves the inequality. At $d=0$ the difference is identically zero.
This argument acts on the continuous payoff as a pointwise function and
does not differentiate a finite atomic probability law.
\end{proof}

\subsection{Matched finite raw-payoff implementation}

Retain $K=100$, $H=1$, $q=28$ and the guarded graph. Conditional on an
actual outer coordinate $y$, its decoded endpoints satisfy exactly
\begin{equation}
 \widehat X_G=\widehat m+\lambda\widetilde N+\zeta(\widetilde N),\qquad
 \widehat X_A=\widehat X_G+\widehat\Delta,\qquad
 |\widehat m-m|\leq E_m=e_m+e_C+e_O,\quad
 \widehat\Delta\geq0.
 \label{eq:common-information-raw-endpoints}
\end{equation}
Here $C=J-B_G$ and $O=J+\lambda L$ are constants, and their errors
$e_C,e_O$ are independent of the inner label. The exact emitted scaling
coefficient is $\widehat\lambda=I_\lambda2^{-d_\lambda}$ and
\begin{equation}
 |\zeta|\leq E_s=2L|\widehat\lambda-\lambda|
                    +2^{-q}\mathbf1_{\{d_\lambda>0\}}.
 \label{eq:common-information-scale-error}
\end{equation}
Let the two distinct scalar maps have intrinsic exponential errors
$e_{\exp,G},e_{\exp,A}$ on their guarded domains. Intrinsic errors exclude
propagated argument perturbations. Let the rounded inner normal law have
both right-CDF and left-CDF discrepancy at most $D_N$ from $\Phi$.
The prepared flag encodes the exact projected dyadic payoff difference.
Its conditional probability $p(y)$ obeys
\begin{equation}
 |p(y)-R_1(y)|\leq\eta_{\rm raw}:=
 \min\left\{1,100(e_{\exp,G}+e_{\exp,A})+101E_s+2D_N
                       +C(\Delta_*)E_m+c_0e_\Delta\right\}.
 \label{eq:common-information-raw-bound}
\end{equation}
An independently valid pointwise or earlier smoothing bound may be
included in this minimum.

To prove \eqref{eq:common-information-raw-bound}, replace the two
represented scalar maps by exact payoffs at their represented arguments.
Their separate errors contribute $100(e_{\exp,G}+e_{\exp,A})$.
Lemma~\ref{lem:common-information-inner-payoff} removes the common
label-dependent $\zeta$ for a charge $101E_s$. Projection is
$1$-Lipschitz and fixes the resulting ideal nonnegative difference.
That difference as a function of the normal coordinate has bounded
variation at most two; integration by parts against the signed CDF
difference therefore charges at most $2D_N$. This statement retains
atomic label multiplicities and uses the left limit at strict endpoints.
Only after this finite-to-continuous transfer is the Gaussian
convolution differentiated. Theorem~\ref{thm:common-information-stability}
then charges $C(\Delta_*)E_m+c_0e_\Delta$. The total is capped at one
because both target and prepared probability belong to $[0,1]$.
Normal quantile rounding is contained in $D_N$; it is not also included
in $E_s$ or $E_m$. Outer quantile rounding remains solely in the existing
financial outer-law transfer.

\subsection{Matched finite numerical conditional-price implementation}

For the retained numerical pool, the cap and the positive monetary
normaliser are exactly $H=B_1=1$. Consequently its two scaling nodes
multiply by the exact integer one and introduce neither coefficient
rounding nor a fixed-point floor. Write the scalar domains as
$[-B_G,B_G]$ and $[-B_A,B_A]$. The encoded nodes satisfy
\begin{equation}
 \widehat M_A=\widehat M_G+\widehat C_{AG}+\widehat\Delta,\qquad
 C_{AG}=B_A-B_G,\qquad
 \widehat m=\widehat M_G-B_G,\qquad
 \widehat d=\widehat\Delta+(\widehat C_{AG}-C_{AG}).
 \label{eq:common-information-numerical-endpoints}
\end{equation}
Suppose the retained actual-range certificate additionally shows that
the input clamps preceding both price tables are inactive on every
actual guarded input. A function-input domain check after a clamp alone
does not establish this premise. Let $e_{P,G},e_{P,A}$ be the intrinsic
conditional-price table errors, $e_{AG}$ the constant error and $e_m$
the lower encoded node error. The prepared projected difference then
satisfies
\begin{equation}
 |p_{\rm num}(y)-R_1(y)|\leq\eta_{\rm num}:=
 \min\{1,e_{P,G}+e_{P,A}+C(\Delta_*)e_m
                             +c_0(e_\Delta+e_{AG})\}.
 \label{eq:common-information-numerical-bound}
\end{equation}
Indeed, replacing both tables by their exact scalar functions at the
actual arguments costs their summed intrinsic errors. Exact decoded
addition gives \eqref{eq:common-information-numerical-endpoints}.
The effective gap need not be nonnegative, but the path bound in
Theorem~\ref{thm:common-information-stability} remains valid and the
final signed projection fixes the true correction. In the retained
graphs $B_A-B_G$ is dyadic and exactly encoded; retaining $e_{AG}$ in
the rule makes the required premise explicit. There is no sampled
inner Gaussian or strike law in this implementation, so adding $D_N$
or $\beta_S$ would charge an approximation it does not perform.
Normalisers other than one require separately certified scaling
coefficient and floor allowances and are not covered by this specialised
numerical formula.

\subsection{Matched finite shared-threshold implementation}

For a fixed ideal uniform strike midpoint $s$, put $a(s)=\log(1+s/100)$.
The exact graph computes a lower encoded endpoint, its upper endpoint
and one shared threshold as
\begin{equation}
 \widehat\ell=\widehat M_G+\widehat C,\qquad
 \widehat u=\widehat\ell+\widehat\Delta,\qquad
 \widehat z=\widehat a(s)+\widehat O
                       -\mathsf S_\lambda(\widetilde N+L).
 \label{eq:common-information-threshold-endpoints}
\end{equation}
Its predicate is $\widehat\ell<\widehat z\leq\widehat u$ and the gap
word is nonnegative. The represented strike error
$|\widehat a(s)-a(s)|\leq e_a$ includes its input-ratio arithmetic
and intrinsic logarithm error; it does not include normal quantile
rounding. The same $C,O,E_s$ as above identify the conceptual threshold
$a(s)-\lambda\widetilde N-\zeta(\widetilde N)$ relative to a common
represented lower endpoint. The represented strike and common endpoint
errors give a fixed-strike common shift bounded by
\begin{equation}
 E_{m,a}=e_m+e_C+e_O+e_a.
 \label{eq:common-information-threshold-common-error}
\end{equation}
They may depend arbitrarily on the fixed outer and strike labels.
They are independent of the inner normal label, which is the required
conditioning condition. With the retained strike midpoint bias $\beta_S$
in probability units, the implemented conditional event probability
obeys
\begin{equation}
 |p_{\rm thr}(y)-R_1(y)|\leq\eta_{\rm thr}:=
 \min\{1,\beta_S+2D_N+2c_0E_s
                     +C(\Delta_*)E_{m,a}+c_0e_\Delta\}.
 \label{eq:common-information-threshold-bound}
\end{equation}

For completeness, at a fixed strike and outer coordinate the bounded
$\zeta$ moves each of the two strict/non-strict Gaussian boundaries by
at most $E_s$. The represented event is contained in the interval
obtained by expanding the unperturbed interval by $E_s$ at each end.
It contains the interval obtained by shrinking by $E_s$ at each end,
whenever that interval is nonempty. The probability of either fixed
interval under the finite normal law differs from the continuous normal
probability by at most $2D_N$: an interval probability is a difference
of two right-CDF or left-CDF values. The continuous density is bounded
by $c_0$, so expansion or shrinkage changes probability by at most
$2c_0E_s$. If the shrunken interval is empty, its original width is
at most $2E_s$ and its entire continuous probability is bounded by
$2c_0E_s$. Thus the finite event transfer costs $2D_N+2c_0E_s$ once.
After this transfer, each fixed-strike continuous interval is the
Gaussian correction with threshold measure $\delta_{a(s)}$, which has
mass one. The common shift and gap allowances follow from
Theorem~\ref{thm:common-information-stability}. Averaging those uniform
allowances over the equal-weight strike labels preserves them; replacing
that strike law by the ideal uniform law charges the separately retained
$\beta_S$. This proves \eqref{eq:common-information-threshold-bound}.
Finite atoms are retained, and no second inner Gaussian-box mass is
multiplied into the decoder.

The label-dependent $\zeta$ cannot be charged merely with
$C(\Delta_*)E_s$ by invoking Gaussian convolution: it is correlated with
the normal label and the conditioning premise of that invocation fails.
The preceding sandwich proves the required allowance without such an
independence assumption.

\subsection{Gap-sensitive finite-law allowances and their limits}

The exact capped exponential difference is unimodal. Its total variation
in the normal coordinate is
\begin{equation}
 \operatorname{TV}\{f(m+\Delta+\lambda\,\cdot)-f(m+\lambda\,\cdot)\}
 =2v(\Delta),\qquad
 v(\Delta)=\min\{1,101(1-e^{-\Delta})\}.
 \label{eq:common-information-gap-variation}
\end{equation}
The retained earlier piecewise derivation proves this exact expression.
Accordingly, the raw finite-law transfer can use
$2D_Nv(\widehat\Delta)$; a true gap bound alone must first be enlarged
by its represented gap error. If $U=e^\Delta-1\geq0$, then
$v(\Delta)=\min\{1,101U/(1+U)\}$. This is a nondecreasing concave
function of $U$. Given an independently certified finite outer moment
$\mathbb E U\leq\overline U$, Jensen's inequality and the Lipschitz
constant $101$ of $v$ yield
\begin{equation}
 \mathbb E v(\widehat\Delta)
 \leq\min\left\{1,
       \min\left\{1,\frac{101\overline U}{1+\overline U}\right\}
                   +101e_\Delta\right\}.
 \label{eq:common-information-gap-mean}
\end{equation}
This is an averaged outer-law certificate, not a uniform conditional
certificate or an automatic two-sided prepared-probability prior.
For the four retained winning bases the certified first moments are
approximately $0.10888$, $0.02563$, $0.10684$ and $0.10666$.
The inner minimum in \eqref{eq:common-information-gap-mean} equals one
in all four cases. This moment-only refinement therefore provides no
positive cost consequence in the retained study. A tighter polynomial
moment majorant would require its own proved domination, optimisation
record and complete prior/budget reconstruction.

All three prepared methods can use their corresponding common-information
formula with the same $\nu_0,\nu_1$, true gap premise and financial
information. For the implemented basket, $\Delta=\log(1+U)\leq U\leq U_*$
provides a price-independent rational choice $\Delta_*=U_*$. A smaller
certified gap bound may be used only when retained with its derivation.
Any smaller probability allowance must be propagated through the
unchanged analytic coefficient, external outer-law mass, continuous
financial transfer, finite mean prior, estimation accuracy, confidence
and complete clean circuit cost. A certificate improvement changes
neither the ideal financial price nor the generic amplitude-estimation
query law.
